\documentclass[10pt,a4paper]{amsart}
\usepackage[margin=19mm]{geometry}
\usepackage{amsmath,amssymb,mathtools}
\usepackage[scr=rsfs,cal=euler]{mathalfa}
\usepackage{xfrac}
\usepackage[unicode,hidelinks,bookmarks=false]{hyperref}
\newcommand{\ie}{i.e.\ }
\newcommand{\cf}{cf.\ }
\newcommand{\resp}{respectively\ }
\newcommand{\Subsection}{\subsection}
\providecommand{\nobreakdash}{\nobreak\hbox{-}}
\setlength{\emergencystretch}{2em}
\multlinegap0pt
\usepackage{amssymb}
\usepackage{xcolor}
\usepackage{float}
\usepackage{tikz}
\newtheorem{proposition}{Proposition}
\newcommand{\hC}{\widehat{C}}
\DeclareMathOperator{\proj}{P}
\title{Split Casimir Operator of the Lie Algebra so(2r) in Spinor Representations, Colour Factors and Yang-Baxter Equation}
\author{A. P. Isaev (1, 2), A. A. Provorov (1) ((1) Bogoliubov Laboratory of Theoretical Physics, JINR, Dubna, (2) Lomonosov Moscow State University, Physics Faculty)}
\date{Source: arXiv:2603.04561v1 (CC BY 4.0)}
\begin{document}
\maketitle

\noindent\emph{Source and licence.} Split Casimir Operator of the Lie Algebra so(2r) in Spinor Representations, Colour Factors and Yang-Baxter Equation. Version arXiv:2603.04561v1 (CC BY 4.0), distributed by arXiv under the Creative Commons Attribution 4.0 International licence, http://creativecommons.org/licenses/by/4.0/, as recorded in the arXiv licence metadata for this exact version and shown on the abstract page https://arxiv.org/abs/2603.04561v1. Full licence notice: \texttt{licenses/arxiv-2603.04561-CC-BY-4.0.txt}.\par\smallskip

\noindent\textit{Editorial scope.} Counted unit: the article's proposition idens (source0.tex lines 599--604). The article's complete original proof of it is delivered with it (source0.tex lines 605--623, one \texttt{proof} environment of 19 lines). No mathematics is written, repaired or re-derived: the statement and the proof are the article's own lines, in the article's own notation, and no retained line is substituted or re-flowed.  Retained uncounted as the source-local context the proof needs: lines 416--419; lines 434--438; lines 557--562; lines 1817--1819.  Nothing was omitted from any retained span, so the package carries no omission record.  No retained line contains a citation, so the wrapper carries no bibliography and no bibliography entry.  No retained line contains a figure environment, an image-inclusion command or drawing code, so no image is delivered and the package declares no asset dependency.  Editorial formatting only, in addition: an AMS \texttt{amsart} preamble that restates the article's own declarations and macros used by the retained lines, namely the environments \texttt{proposition} on separate counters, together with the standard packages those lines need.  The retained environments print in this document as Proposition 1 in order of appearance, whereas the article prints the same items with the numbering of its own full document.  The complete native source of the article as distributed in the arXiv e-print for this exact version is bundled unchanged as \texttt{sources/195748/source0.tex}.
\begin{equation}\label{I2 via hC}
I_2=\Gamma^{[i_1 i_2]}\otimes \Gamma_{[i_1 i_2]}
=-16(N-2)\hC_\rho,
\end{equation}

\begin{equation}\label{I2kI2 relation}
I_{2k} I_2 = I_{2k+2}
+ 8k(r-k) I_{2k}
+ 4k(2k-1)(r+1-k)(2r+1-2k) I_{2k-2}.
\end{equation}

\begin{equation}\label{(I otimes Gamma)I_k}
(I\otimes \Gamma_{2r+1})\frac{I_k}{k!}
=
(-1)^r(\Gamma_{2r+1}\otimes I)
\frac{I_{2r-k}}{(2r-k)!}.
\end{equation}

\begin{equation}\label{Gamma_N+1 properties}
\Gamma_{2r+1}^2=I, \qquad \Gamma_{2r+1}\Gamma_i=-\Gamma_i\Gamma_{2r+1} \quad \text{for} \quad i=1,\dots,2r.
\end{equation}

\begin{proposition}
The polynomials $I_{2k}(\hC_{\epsilon\epsilon'})$ satisfy
\begin{equation}\label{I_k's: hC possible idens}
I_{2r-2k}(\hC_{\epsilon\epsilon'})-\epsilon\epsilon' \frac{(2r-2k)!}{(2k)!}I_{2k}(\hC_{\epsilon\epsilon'})=0,\quad k=0,1,\dots,2r.
\end{equation}
\end{proposition}

\begin{proof}
For any $k=0,1,\dots,2r$, consider the following chain of equalities:
\begin{equation}\label{proj_ee' I_k}
\begin{aligned}
\proj_{\epsilon\epsilon'}\frac{I_k}{k!}&=\frac{1}{4}\big((I+\epsilon \Gamma_{{2r}+1})\otimes (I+\epsilon' \Gamma_{{2r}+1})\big)\frac{I_k}{k!}\\
&=\frac{1}{4k!}(I\otimes I)I_k+\frac{1}{4k!}\epsilon(\Gamma_{{2r}+1}\otimes I)I_k+\frac{1}{4k!}\epsilon'(I\otimes \Gamma_{{2r}+1})I_k+\frac{1}{4k!}\epsilon\epsilon'(\Gamma_{{2r}+1}\otimes \Gamma_{{2r}+1})I_k\\
&=\frac{(-1)^r}{4({2r}-k)!}(\Gamma_{{2r}+1}\otimes \Gamma_{{2r}+1})I_{{2r}-k}+\frac{(-1)^r}{4({2r}-k)!}\epsilon (I\otimes \Gamma_{{2r}+1})I_{{2r}-k}+\\
&+\frac{(-1)^r}{4({2r}-k)!}\epsilon'(\Gamma_{{2r}+1}\otimes I)I_{{2r}-k}+\frac{(-1)^r}{4({2r}-k)!}\epsilon\epsilon' (I\otimes I)I_{{2r}-k}\\
&=(-1)^r\frac{1}{4}\big((\epsilon I+\Gamma_{{2r}+1})\otimes (\epsilon' I+\Gamma_{{2r}+1})\big)\frac{I_{{2r}-k}}{({2r}-k)!}=\epsilon\epsilon' (-1)^r\proj_{\epsilon\epsilon'}\frac{I_{{2r}-k}}{({2r}-k)!}.
\end{aligned}
\end{equation}
Here we have used the involutive property~\eqref{Gamma_N+1 properties}
of $\Gamma_{2r+1}$ and the multiplication rule~\eqref{(I otimes Gamma)I_k}
for $(I\otimes \Gamma_{2r+1})$ and $(\Gamma_{2r+1}\otimes I)$ acting on $I_k$.
Then~\eqref{I_k's: hC possible idens} follows from~\eqref{proj_ee' I_k} by substituting $k\mapsto 2k$ and expressing the invariants
$\proj_{\epsilon\epsilon'} I_{2k}$ and $\proj_{\epsilon\epsilon'} I_{2r-2k}$
as polynomials in $\hC_{\epsilon\epsilon'}$
using~\eqref{I2 via hC} and~\eqref{I2kI2 relation}.
\end{proof}

\end{document}
